\begin{helper}
Under the two-bite law, the probability of the projection-size event \(\noderef{TwoBiteProjectionSizeEvent}(I,k,\ell_R,\ell_B)\) is at most
\[
\frac{\binom{m}{\ell_R}\binom{m}{\ell_B}\binom{\ell_R\ell_B}{k}}{\binom{m^2}{k}}.
\]
\end{helper}
\begin{proof}
Let \(U=\operatorname{Fin}m\times\operatorname{Fin}m\), and let
\(\Omega=\{\pi:\operatorname{Fin}n\hookrightarrow U\}\).  For an injection
\(\pi\), put \(S_\pi=\pi(I)\).  Since \(|I|=k\), injectivity gives
\(|S_\pi|=k\).  The event
\(\noderef{TwoBiteProjectionSizeEvent}(I,k,\ell_R,\ell_B)\) depends only on
\(\pi\): by the definitions of the red and blue projections,
\[
\pi_R(I)=\operatorname{fst}(S_\pi),\qquad
\pi_B(I)=\operatorname{snd}(S_\pi).
\]
Thus the event is exactly the injection-only property
\[
P(S_\pi):\quad
|S_\pi|=k,\quad |\operatorname{fst}(S_\pi)|=\ell_R,\quad
|\operatorname{snd}(S_\pi)|=\ell_B.
\]

By \(\noderef{TwoBiteEventProbabilityInjectionOnly}\), the two-bite
probability of this event is
\[
N_{\rm inj}\cdot \noderef{UniformInjectionWeight}(n,m),
\]
where \(N_{\rm inj}\) is the number of injections \(\pi\in\Omega\) for which
the above property holds.

We first dispose of zero cases.  If \(\binom{m^2}{k}=0\), then there is no
\(k\)-subset of \(U\).  Since every injection would produce the
\(k\)-element subset \(S_\pi\), there are no relevant injections; the
probability is \(0\).  In the Lean real-number statement, division by this
zero denominator is interpreted as \(0\), so the displayed bound also has
right hand side \(0\).  If \(|\Omega|=0\) while \(\binom{m^2}{k}>0\), then
\(\noderef{UniformInjectionWeight}(n,m)=0\), so the probability is \(0\),
whereas the right hand side is nonnegative.  The inequality follows in both
zero cases.

Assume from now on that \(|\Omega|>0\) and \(\binom{m^2}{k}>0\).  Then
\[
\noderef{UniformInjectionWeight}(n,m)=\frac1{|\Omega|},
\]
so the probability is \(N_{\rm inj}/|\Omega|\).

Apply \(\noderef{UniformInjectionImage}\) to the property
\[
Q(S):\quad |\operatorname{fst}(S)|=\ell_R
\text{ and }|\operatorname{snd}(S)|=\ell_B.
\]
For every injection \(\pi\), the set \(S_\pi\) already has size \(k\), so
\(Q(S_\pi)\) is equivalent to
\(\noderef{TwoBiteProjectionSizeEvent}(I,k,\ell_R,\ell_B)\).  Therefore
\[
N_{\rm inj}\binom{m^2}{k}
=N_{\rm sub}|\Omega|,
\]
where \(N_{\rm sub}\) is the number of \(k\)-subsets
\(S\subseteq U\) with
\[
|\operatorname{fst}(S)|=\ell_R,\qquad
|\operatorname{snd}(S)|=\ell_B.
\]
Because both denominators are positive, we may divide by
\(|\Omega|\binom{m^2}{k}\) and obtain
\[
\frac{N_{\rm inj}}{|\Omega|}
=\frac{N_{\rm sub}}{\binom{m^2}{k}}.
\]

Finally, \(\noderef{ProjectionSubsetCountBound}\) gives
\[
N_{\rm sub}\le
\binom{m}{\ell_R}\binom{m}{\ell_B}\binom{\ell_R\ell_B}{k}.
\]
Dividing by the positive \(\binom{m^2}{k}\) yields
\[
\operatorname{TwoBiteEventProbability}
(\noderef{TwoBiteProjectionSizeEvent}(I,k,\ell_R,\ell_B))
\le
\frac{\binom{m}{\ell_R}\binom{m}{\ell_B}\binom{\ell_R\ell_B}{k}}
{\binom{m^2}{k}},
\]
as required.
\end{proof}
